\documentclass[11pt]{article}

\usepackage[T1]{fontenc}
\usepackage[utf8]{inputenc}
\usepackage[margin=1in]{geometry}
\usepackage{amsmath, amssymb, amsthm}
\usepackage{hyperref}

\hypersetup{
  colorlinks=true,
  linkcolor=blue,
  urlcolor=blue,
  citecolor=blue
}

\title{Architecture Proof Notes for Modern LLM Families}
\author{Fractal Research Notes}
\date{\today}

\newtheorem{definition}{Definition}[section]
\newtheorem{lemma}[definition]{Lemma}
\newtheorem{theorem}[definition]{Theorem}
\newtheorem{proposition}[definition]{Proposition}
\newtheorem{conjecture}[definition]{Conjecture}
\newtheorem{remark}[definition]{Remark}

\newcommand{\Softmax}{\operatorname{Softmax}}
\newcommand{\Attn}{\operatorname{Attn}}
\newcommand{\diag}{\operatorname{diag}}
\newcommand{\R}{\mathbb{R}}

\begin{document}

\maketitle

\begin{abstract}
These notes collect the smallest clean theorem surface for four architecture
families:
\begin{enumerate}
  \item standard decoder-only language models,
  \item modern sub-variants including sparse mixture-of-experts, linear
    attention, and hybrid recurrent-attention stacks such as Kimi Linear,
  \item the current \texttt{A + P2} contender family,
  \item a speculative Graph-of-Experts / thought-channel family for bounded
    internal search.
\end{enumerate}

The aim is not to prove empirical superiority.
The aim is to isolate claims that are actually provable:
causality, incremental equivalence, reduction to simpler cases, strict
structural generalization where valid, and asymptotic or control-plane bounds.
\end{abstract}

\section{Proof Discipline and Common Machine}

\begin{remark}
Every formal statement in this note should be interpreted as exactly one of:
\emph{definition}, \emph{lemma}, \emph{theorem}, \emph{proposition}, or
\emph{conjecture}. If benchmark evidence is required, the statement is not a
theorem.
\end{remark}

Let the input token sequence be
\[
  x_1, x_2, \dots, x_T.
\]
At step $t$, define:
\begin{itemize}
  \item $h_t$: active latent state,
  \item $M_t$: memory object,
  \item $c_t$: control object,
  \item $y_t$: emitted token representation used by the output head.
\end{itemize}

We place all families inside the generic causal machine
\begin{align}
  (h_t, M_t, c_t) &= U_\theta(h_{t-1}, M_{t-1}, c_{t-1}, x_t), \\
  y_t &= R_\theta(h_t, M_t, c_t, x_t), \\
  p_\theta(x_{t+1}\mid x_{\le t}) &= \Softmax(W y_t).
\end{align}

What varies across architectures is:
\begin{itemize}
  \item what $M_t$ stores,
  \item what $c_t$ controls,
  \item whether $U_\theta$ is attention-dominant, recurrent-dominant, hybrid,
    or search-native,
  \item how much exact token interaction is available inside $U_\theta$.
\end{itemize}

\subsection{Shared Symbols}

\begin{itemize}
  \item $T$: sequence length,
  \item $d$: model width,
  \item $w$: local exact-attention window,
  \item $L$: total layer count,
  \item $L_A$: exact-attention layer count,
  \item $L_S$: sequence-mixer layer count,
  \item $E$: number of experts in a sparse MoE layer,
  \item $k$: number of active experts per token,
  \item $K$: number of active thought channels,
  \item $B$: bounded internal search budget,
  \item $s_t$: recurrent latent state in a sequence-mixer block.
\end{itemize}

\subsection{Basic Causality}

\begin{definition}[Causal update]
An update rule is \emph{causal} if step $t$ depends only on $x_{\le t}$ and
state produced from steps $< t$, and never on any future token $x_u$ with
$u > t$.
\end{definition}

\begin{lemma}[Causal composition]
If $f$ and $g$ are causal update operators over the same prefix order, then
their composition $g \circ f$ is causal.
\end{lemma}

\begin{proof}
The output of $f$ depends only on allowed prefix information. The output of
$g$ depends only on the output of $f$ and the same allowed prefix information.
Therefore no future token can enter through composition.
\end{proof}

\begin{definition}[Incremental-decode equivalence]
An incremental decode procedure is equivalent to full-prefix recomputation if,
for fixed weights, masks, and prefix,
\[
  y_t^{\mathrm{inc}} = y_t^{\mathrm{full}}
\]
up to deterministic arithmetic tolerance.
\end{definition}

\begin{definition}[Reduction to a simpler class]
Architecture family $F$ reduces to family $G$ if there exists an explicit
parameter or control setting under which every member of the reduced subclass
of $F$ behaves as a member of $G$.
\end{definition}

\begin{definition}[Strict structural generalization]
Family $F$ strictly structurally generalizes family $G$ if:
\begin{enumerate}
  \item $F$ reduces to $G$ under an explicit restriction, and
  \item there exist members of $F$ that cannot be represented inside $G$
    without changing the declared contract class.
\end{enumerate}
\end{definition}

\section{Standard Decoder LLMs}

Assume masked causal self-attention, residual blocks, and position-wise
feed-forward updates.

For one attention head:
\begin{align}
  q_t &= W_Q r_t, \\
  k_u &= W_K r_u, \\
  v_u &= W_V r_u.
\end{align}
With causal masking:
\[
  \alpha_{t,u} =
  \begin{cases}
    0, & u > t, \\
    \dfrac{\exp(q_t \cdot k_u / \sqrt{d_k})}
      {\sum_{j \le t}\exp(q_t \cdot k_j / \sqrt{d_k})}, & u \le t.
  \end{cases}
\]
Then
\[
  \Attn_t = \sum_{u \le t} \alpha_{t,u} v_u.
\]

\begin{definition}[Decoder language model]
A standard decoder language model is any stacked residual architecture in which
every layer is causal and the next-token distribution is read from the top
layer token representation:
\[
  p_\theta(x_{t+1}\mid x_{\le t}) = \Softmax(W r_t^{(L)}).
\]
\end{definition}

\begin{theorem}[Causal decoder correctness]
A masked decoder transformer defines a causal conditional distribution over the
next token.
\end{theorem}

\begin{proof}
At every layer, the causal mask forbids access to positions $u > t$.
Feed-forward blocks are pointwise in position and therefore cannot introduce
future-token dependence. Residual addition preserves causality whenever both
branches are causal. By repeated application of the causal-composition lemma,
the full stack remains causal, and the output head reads only the final causal
state at position $t$.
\end{proof}

\begin{theorem}[KV-cache decode equivalence]
Under fixed weights, fixed masks, and deterministic inference,
incremental decode with cached keys and values is equivalent to full-prefix
recomputation for the same prefix.
\end{theorem}

\begin{proof}
At step $t$, full-prefix recomputation attends to the set
$\{(k_u, v_u) : u \le t\}$. Under fixed weights, the keys and values for prior
positions $u < t$ do not change when a new token arrives. Therefore the cached
set of prior key-value pairs is exactly the same set that full-prefix
recomputation would rebuild. The only new objects required are the current
token's key and value. Since the score function and mask are unchanged, the
attention output is the same.
\end{proof}

\begin{proposition}[Approximate copy capacity]
Exact self-attention can implement arbitrarily sharp copy-like retrieval of a
visible prior token.
\end{proposition}

\begin{proof}
Choose parameters so that the query at position $t$ scores one visible prior
position $u^\star \le t$ strictly above all others. Then the attention weights
concentrate on $u^\star$, and the resulting output becomes arbitrarily close to
$v_{u^\star}$. If the output projection reads the copied content from
$v_{u^\star}$, the model implements pointer-style copying over the visible
context.
\end{proof}

\begin{proposition}[Token comparison capacity]
Exact self-attention can implement local or long-range token comparison over
the visible context.
\end{proposition}

\begin{proof}
Keys and queries can be parameterized so that the current token representation
selects one or more comparison-relevant prior tokens. The attention read gathers
their information into the residual stream, and a subsequent feed-forward block
can map that gathered evidence to a comparison feature.
\end{proof}

\begin{proposition}[Complexity bound]
Let sequence length be $T$, model width be $d$, and layer count be $L$. Then
train-time dense attention is quadratic in visible sequence length per
attention layer, while decode-time incremental attention is linear in visible
cache length per attention layer.
\end{proposition}

\begin{proof}
At train time, each visible query position scores against all visible prior
keys in each attention layer. That induces the familiar quadratic prefix
interaction. At incremental decode time, each new query scores against the
already-visible cache once per layer, which is linear in cache length for that
step.
\end{proof}

\section{Modern Sub-Variants}

\subsection{Sparse Mixture of Experts}

Write one sparse-MoE feed-forward block as
\[
  \mathrm{MoE}(x_t) = \sum_{i \in \mathrm{Active}_t} \pi_{t,i} E_i(x_t),
\]
where $|\mathrm{Active}_t| \le k$.

\begin{proposition}[Dense reduction for sparse MoE]
A sparse-MoE layer reduces to a dense feed-forward layer under a degenerate
expert-sharing and routing choice.
\end{proposition}

\begin{proof}
Set every expert function equal to the same dense function $F$. Force the
router to activate exactly one expert with weight $1$. Then
\[
  \mathrm{MoE}(x_t) = F(x_t).
\]
Hence dense FFN is a special case of MoE.
\end{proof}

\begin{proposition}[Expert budget invariant]
If routing activates at most $k$ experts per token, then per-token expert
compute is bounded by $k$ expert evaluations.
\end{proposition}

\begin{proof}
The router is the only component allowed to activate experts. By contract it
emits at most $k$ active experts for a token, so no more than $k$ expert
evaluations can occur for that token.
\end{proof}

\begin{proposition}[Control-plane separation]
Sparse MoE changes conditional compute allocation without changing the causal
prefix contract of the surrounding decoder.
\end{proposition}

\begin{proof}
Routing depends on the current causal token state, and expert evaluation
consumes only that causal state. Therefore expert selection does not introduce
future-token dependence by itself.
\end{proof}

\begin{conjecture}[Parameter-efficiency benefit]
Sparse expert specialization may improve parameter-efficiency. This is an
empirical claim, not a theorem.
\end{conjecture}

\subsection{Linear-Attention Families}

For a factorized linear-attention family with feature map $\phi$, define
\begin{align}
  S_t &= \sum_{u \le t} \phi(k_u) v_u^\top, \\
  z_t &= \sum_{u \le t} \phi(k_u), \\
  \mathrm{output}_t &= \frac{\phi(q_t)^\top S_t}{\phi(q_t)^\top z_t}.
\end{align}

\begin{theorem}[Associative recurrent form]
If the attention kernel factors through an associative feature map $\phi$, then
prefix processing admits an equivalent recurrent update with linear-time
sequence evolution.
\end{theorem}

\begin{proof}
The objects $S_t$ and $z_t$ can be updated from $S_{t-1}$ and $z_{t-1}$ using
only the current token contribution. No revisit of all prior tokens is required
once those prefix summaries are maintained. Therefore the visible-prefix read
admits a recurrent implementation over bounded summary state.
\end{proof}

\begin{remark}[Softmax exactness boundary]
The factorized recurrent form above should not be read as a general
exact-equivalence guarantee to arbitrary softmax attention.
Establishing a sharp non-equivalence theorem would require a separate
representation argument beyond the scope of these notes.
\end{remark}

\begin{proposition}[Decode complexity bound]
If the model maintains recurrent prefix summaries of bounded size, incremental
decode cost is bounded by the recurrent-state update rather than a full
visible-cache scan.
\end{proposition}

\begin{proof}
The recurrent summaries replace explicit re-scanning of all visible prior
tokens. Thus the decode cost is governed by summary update size rather than
visible context length itself.
\end{proof}

\subsection{Hybrid Recurrent-Attention Families}

This class includes:
\begin{itemize}
  \item Jamba-style interleaved attention and SSM hybrids,
  \item Nemotron-H-style mostly-SSM stacks with occasional attention,
  \item Kimi-style recurrent-memory plus periodic exact-attention stacks.
\end{itemize}

Let the stack contain $L_A$ exact-attention layers and $L_S$ recurrent or SSM
sequence-mixer layers.

\begin{definition}[Hybrid layer cost model]
Let
\[
  c_A(T, d), \quad m_A(T, d)
\]
denote the per-layer compute and memory obligations of one exact-attention
layer on visible context length $T$, and let
\[
  c_S(d), \quad m_S(d)
\]
denote the per-layer compute and memory obligations of one recurrent or SSM
layer with bounded summary state.
\end{definition}

\begin{theorem}[Hybrid causal composition]
An interleaved stack of causal exact-attention and causal recurrent or SSM
blocks remains causal.
\end{theorem}

\begin{proof}
Each attention block is causal by masked attention. Each recurrent or SSM block
is causal by left-to-right update. The full stack is causal by repeated
composition.
\end{proof}

\begin{proposition}[Exact-attention budget decomposition]
Under the hybrid layer cost model, the exact-attention-specific portion of the
stack cost is
\[
  L_A \, c_A(T, d)
\]
in compute and
\[
  L_A \, m_A(T, d)
\]
in memory.
\end{proposition}

\begin{proof}
Only the exact-attention layers incur exact visible-token interaction costs.
The recurrent or SSM layers replace those exact-attention obligations with
their own bounded summary-state updates. Summing over $L_A$ exact-attention
layers yields the stated exact-attention compute and memory terms.
\end{proof}

\begin{proposition}[Reduction edge cases]
A hybrid recurrent-attention stack reduces to attention-only when $L_S = 0$,
and to recurrent-only when $L_A = 0$.
\end{proposition}

\begin{proof}
Setting one layer family count to zero removes that family entirely. The
remaining stack is exactly the edge-case class.
\end{proof}

\begin{proposition}[Kimi-style periodic exact-attention exposure]
In a fixed $3{:}1$ KDA-to-MLA schedule, an exact-attention layer occurs at
regular depth intervals while recurrent memory carries the cheaper background
state between those refresh points.
\end{proposition}

\begin{proof}
Every fourth layer is an exact-attention layer in the stated schedule.
Therefore any representation that traverses depth encounters an exact
interaction surface at regular intervals, while the intervening KDA layers
carry the recurrent-memory path.
\end{proof}

\begin{proposition}[Hybrid memory decomposition]
Under the hybrid layer cost model, total layer-local memory obligation is
\[
  L_A \, m_A(T, d) + L_S \, m_S(d).
\]
In particular, the exact visible-token memory term scales with $L_A$ rather
than the full layer count $L$.
\end{proposition}

\begin{proof}
The recurrent-memory layers replace full exact-attention visible-token state
with bounded recurrent summaries. Therefore exact visible-token memory appears
only in the $L_A$ exact-attention layers, while the remaining $L_S$ layers
contribute only their bounded summary-state terms.
\end{proof}

\section{\texttt{A + P2}}

The \texttt{A + P2} family is the current Path 1 contender:
\[
  A\text{-}P2\text{-}A\text{-}P2\text{-}A\text{-}P2\text{-}A\text{-}P2.
\]

\subsection{Expected Contract Class}

We assume the currently expected direction:
\begin{itemize}
  \item predictive-core sequence primitive only,
  \item no memory/index sidecar behavior,
  \item transformed state dynamics before blending,
  \item emitted output distinct from latent state,
  \item no routed retrieval or Path 2 control-plane behavior.
\end{itemize}

Write the \texttt{P2} step abstractly as
\begin{align}
  u_t &= T_\theta(s_{t-1}, x_t), \\
  c_t &= C_\theta(x_t), \\
  s_t &= U_\theta(u_t, c_t, x_t), \\
  o_t &= R_\theta(s_t, x_t).
\end{align}

\begin{remark}
This section is intentionally conditional on the current expected \texttt{P2}
contract class. It does not claim that every final implementation detail is
already frozen.
\end{remark}

\begin{definition}[Abstract \texttt{P1}-style contract]
For these notes, a \texttt{P1}-style family is any single-state recurrent
predictor of the form
\begin{align}
  s_t &= G_\theta(s_{t-1}, x_t), \\
  y_t &= s_t,
\end{align}
with no separate emitted object beyond the latent state itself.
\end{definition}

\begin{definition}[Step/scan consistency for \texttt{P2}]
A \texttt{P2} implementation is step/scan consistent if repeated
\texttt{step(...)} over a sequence and one \texttt{scan(...)} over the same
sequence produce the same emitted outputs and final latent state, up to
deterministic arithmetic tolerance.
\end{definition}

\begin{theorem}[Causal \texttt{P2} update correctness]
If $T_\theta$, $C_\theta$, $U_\theta$, and $R_\theta$ depend only on $x_t$ and
$s_{t-1}$, then the \texttt{P2} block is causal.
\end{theorem}

\begin{proof}
The transformed carry path sees only the prior latent state and current input.
The candidate path sees only the current input. The update combines only those
causal objects, and the readout sees only the resulting state and current
input. No future token can enter the block.
\end{proof}

\begin{theorem}[\texttt{A + P2} hybrid causality]
If every $A$ block is causal by masked local attention and every \texttt{P2}
block is causal, then the interleaved \texttt{A + P2} stack is causal.
\end{theorem}

\begin{proof}
Apply the causal-composition lemma layer by layer through the interleaved stack.
\end{proof}

\begin{proposition}[\texttt{P1}-style reduction]
Under explicit restrictions, the \texttt{P2} family reduces to a simpler
\texttt{P1}-style single-state direct-readout family.
\end{proposition}

\begin{proof}
Restrict $R_\theta$ to direct state readout so that $o_t = s_t$, absorb
$C_\theta$ into the update map, and choose $T_\theta$ and $U_\theta$ so that
their composition becomes one single-state recurrent map
$G_\theta(s_{t-1}, x_t)$. Under those restrictions the \texttt{P2} step
becomes
\[
  s_t = G_\theta(s_{t-1}, x_t), \qquad y_t = s_t,
\]
which is exactly the declared \texttt{P1}-style contract.
\end{proof}

\begin{proposition}[Structural enlargement over the abstract \texttt{P1} contract]
If the final \texttt{P2} contract permits a separate emitted object
$o_t = R_\theta(s_t, x_t)$ that is not required to equal $s_t$, then the
expected \texttt{P2} class strictly structurally generalizes the abstract
\texttt{P1}-style family defined above.
\end{proposition}

\begin{proof}
The reduction proposition shows that \texttt{P2} contains a restricted
\texttt{P1}-style subclass. But by definition the abstract \texttt{P1}
contract requires $y_t = s_t$ for every step. Any \texttt{P2} instance with
an emitted object $o_t \neq s_t$ on some reachable state trajectory cannot be
represented inside that direct-readout contract without changing the contract
class. Hence the expected \texttt{P2} class strictly enlarges the structure.
\end{proof}

\begin{proposition}[Matched slot complexity]
Assume both \texttt{M} and \texttt{P2} are bounded-state sequence primitives
whose per-layer update costs are independent of full visible context length
beyond the local-attention work already paid by $A$. Then interleaved
\texttt{A + M} and \texttt{A + P2} stacks share the same slot-level asymptotic
form: exact local-attention cost from $A$ plus bounded sequence-primitive cost
from the companion block.
\end{proposition}

\begin{proof}
The $A$ layers retain the exact local-attention obligation. The \texttt{P2}
layers update bounded recurrent or sequence-mixer state instead of performing
full context scans, and the same assumption holds for the reference
\texttt{M} lane. Therefore both hybrids have the same slot-level asymptotic
decomposition: attention work in $A$ and bounded non-attention update work in
the companion sequence primitive.
\end{proof}

\begin{conjecture}[Efficiency-quality advantage]
Whether \texttt{A + P2} matches or beats \texttt{A} or \texttt{A + M} in
quality, memory, or throughput is empirical.
\end{conjecture}

\section{Graph-of-Experts / Thought-Channel Model}

The speculative Graph-of-Experts family contains:
\begin{itemize}
  \item a shared exact-attention trunk,
  \item $K$ active thought channels,
  \item per-channel recurrent or SSM workspaces,
  \item explicit compare, score, prune, merge, and halt control.
\end{itemize}

The goal is not classic sparse expert routing.
The goal is bounded internal multi-hypothesis search.

At step $t$, define:
\begin{itemize}
  \item $h_t$: shared trunk state,
  \item $F_t = \{f_t^{(1)}, \dots, f_t^{(K_t)}\}$: active channel frontier,
  \item $c_t$: controller state.
\end{itemize}

One abstract transition form is:
\begin{align}
  h_t &= \mathrm{SharedAttention}_\theta(h_{t-1}, x_t), \\
  F_t^{\mathrm{seed}} &= \mathrm{Seed}_\theta(h_t, F_{t-1}, c_{t-1}), \\
  F_t^{\mathrm{expand}} &= \mathrm{Expand}_\theta(F_t^{\mathrm{seed}}, h_t), \\
  F_t^{\mathrm{retrieve}} &= \mathrm{Retrieve}_\theta(F_t^{\mathrm{expand}}, h_t), \\
  F_t^{\mathrm{cmp}} &= \mathrm{CompareMerge}_\theta(F_t^{\mathrm{retrieve}}, h_t), \\
  (F_t, c_t) &= \mathrm{Control}_\theta(F_t^{\mathrm{cmp}}, c_{t-1}), \\
  y_t &= \mathrm{Readout}_\theta(h_t, F_t, c_t).
\end{align}

Assume explicit bounds:
\begin{itemize}
  \item $K_t \le K_{\max}$,
  \item at most $B$ internal refinement rounds per output step,
  \item controller operations preserve the channel cap.
\end{itemize}

\begin{definition}[Channel-separated state]
A thought-channel system has channel-separated state if each live hypothesis is
assigned a distinct explicit state slot and any cross-channel interaction occurs
only through declared compare or merge operators.
\end{definition}

\begin{proposition}[Single-channel reduction]
When $K_{\max} = 1$, branching is disabled, and merge and prune are trivial, the
Graph-of-Experts family reduces to a single-stream hybrid predictor.
\end{proposition}

\begin{proof}
With one channel there is no branching, compare and merge become vacuous, and
the controller reduces to a pass-through. The system becomes a shared trunk
plus one recurrent workspace.
\end{proof}

\begin{proposition}[Channel budget invariant]
If the controller keeps at most $K_{\max}$ live channels after every prune
step, then the live channel count remains bounded by $K_{\max}$ at every step.
\end{proposition}

\begin{proof}
The controller is the only component permitted to alter live-channel count. By
contract it outputs no more than $K_{\max}$ channels after each prune step.
Induction over steps yields the invariant.
\end{proof}

\begin{proposition}[Internal search budget bound]
If the controller allows at most $B$ internal refinement rounds per output step
and each round preserves the channel budget invariant, then per-token internal
search compute is bounded by a function of $K_{\max}$, $B$, one shared-trunk
pass, and one compare/control pass per round.
\end{proposition}

\begin{proof}
The shared-trunk pass occurs once. Each refinement round is explicitly capped,
and each round operates under the preserved channel budget invariant. Therefore
the total internal search work is bounded by the declared control-plane budget.
\end{proof}

\begin{proposition}[No architectural forced collapse under separated slots]
If channel state is separated by explicit slots and cross-channel interaction
occurs only through declared compare or merge operators, then the architecture
does not itself force branch collapse between those declared interaction points.
\end{proposition}

\begin{proof}
Each live hypothesis occupies a distinct state slot. Local channel updates act
only on that slot, and cross-channel interaction occurs only at explicit
operators. Therefore immediate branch collapse is not imposed by the
architecture between declared compare or merge steps.
\end{proof}

\begin{proposition}[Bounded beam-style control skeleton]
A Graph-of-Experts controller with explicit seed, expand, score, prune, merge,
and halt operators realizes the same width-bounded successor-selection
structure as a beam-style search loop in latent space.
\end{proposition}

\begin{proof}
Seed initializes the candidate set, expand creates successor states, score
ranks them, prune enforces bounded width, merge collapses compatible states when
desired, and halt terminates the loop. This is exactly the structural control
pattern used by a width-bounded search loop, now expressed over latent channel
states rather than explicit symbolic states.
\end{proof}

\begin{proposition}[Shared-trunk amortization]
Relative to an external GoT harness that re-runs a full model for each branch,
a Graph-of-Experts family with one shared trunk amortizes common token
processing work across all live channels.
\end{proposition}

\begin{proof}
The shared token-context processing occurs once in the trunk. Channel-local
branching happens after that shared computation, so common prefix processing is
not repeated per branch inside the same output step.
\end{proof}

\begin{proposition}[Token-serialization savings]
A latent-channel model can avoid explicit token serialization of intermediate
branches that an external GoT harness would represent as text or extra model
calls.
\end{proposition}

\begin{proof}
Intermediate branch states remain inside the latent frontier $F_t$ rather than
being emitted as visible text tokens. Hence the family has a lower token
overhead floor than an external text-serialized branch harness.
\end{proof}

\begin{conjecture}[Search-efficiency advantage over external GoT]
The Graph-of-Experts family may be more token-efficient or more compute-efficient
than an external GoT harness because of latent branch storage,
shared-trunk amortization, and bounded internal compare/prune loops.
\end{conjecture}

\begin{conjecture}[Native-search advantage over plain hybrids]
The family may outperform single-stream hybrids on tasks that require
maintaining multiple candidate hypotheses over several reasoning steps.
\end{conjecture}

\section{Boundary}

These notes prove or sketch only what can currently be formalized cleanly:
\begin{itemize}
  \item causality,
  \item decode equivalence where appropriate,
  \item reductions,
  \item structural generalization,
  \item explicit budget invariants,
  \item bounded-search expressiveness.
\end{itemize}

They do not prove benchmark wins, training stability, or emergent reasoning
quality. Those remain empirical.

\end{document}
